\section{Compléments sur la courbe de résilience et les indicateurs (Mesurer)}
\label{sec:S3}

\subsection{Paramètres et estimation}
\label{sec:S3.1}

Le tableau~\ref{tab:d6} donne la lecture opérationnelle des sept paramètres de la courbe (équation~(2) du manuscrit).

\begin{table}[htbp]
\caption{Paramètres de la courbe de résilience.}
\label{tab:d6}
\footnotesize
\begin{tabularx}{\textwidth}{@{}P{1.4cm}P{3.4cm}L@{}}
\toprule
\textbf{Paramètre} & \textbf{Signification} & \textbf{Lecture pour la chaîne logistique} \\
\midrule
$k$ & Niveau avant la perturbation & Service nominal du réseau \\
$h$ & Profondeur de la chute & Ampleur de la perte de service au plus fort de la crise \\
$q$ & Écart entre niveau final et niveau initial & Perte durable ($q$ $<$ 0) ou sur-récupération ($q$ $>$ 0) \\
$g$ & Décalage de la chute & Temps pendant lequel le réseau absorbe le choc avant de se dégrader, par ses marges de capacité et, en représentation temporelle, par ses stocks et la matière en route \\
$d$ & Durée entre chute et reprise & Durée du palier dégradé, liée à la persistance du choc, au délai de réaction et, en représentation temporelle, au temps de reconstitution des stocks après le retour des capacités \\
$\alpha$ & Raideur de la chute & Brutalité de la dégradation \\
$\beta$ & Raideur de la reprise & Rapidité du rétablissement \\
\bottomrule
\end{tabularx}
\end{table}

Quelques contraintes garantissent une forme interprétable : le palier dégradé reste positif, le niveau final reste au-dessus du palier, et les raideurs ainsi que la durée du palier ne peuvent descendre sous des valeurs minimales qui produiraient des formes sans signification.

Chaque trajectoire est recalée sur le déclenchement de la perturbation, de sorte que la période précédant la crise fournit directement le niveau nominal. Les sept paramètres sont estimés en minimisant un critère robuste qui pénalise les grands écarts moins fortement que les moindres carrés :

\begin{equation}
L(\theta) = \sum_{i} 2\left(\sqrt{1 + \frac{r_i^2}{f^2}} - 1\right), \qquad r_i = p(x_i;\theta) - \mathrm{IP}(x_i), \qquad f = 0{,}05
\label{eq:S4}
\end{equation}

Ce critère, ses raisons et un exemple d'ajustement sont présentés en section 3.3 du manuscrit (figure~3).

\subsection{Définition et formules des six indicateurs}
\label{sec:S3.2}

Le tableau~\ref{tab:d7} présente les six indicateurs, la dimension de la résilience que chacun décrit et leur sens favorable.

\begin{table}[htbp]
\caption{Les six indicateurs de résilience.}
\label{tab:d7}
\footnotesize
\begin{tabularx}{\textwidth}{@{}P{3.6cm}P{2.3cm}LP{2.1cm}@{}}
\toprule
\textbf{Indicateur} & \textbf{Dimension} & \textbf{Interprétation} & \textbf{Sens favorable} \\
\midrule
IPC, indice de perte cumulée & Perte & Service cumulé perdu pendant la crise & Faible \\
CRD, coefficient de récupération dynamique & Vitesse & Rapidité de la reprise rapportée à la chute et à la durée du palier & Élevé \\
TRP, temps de récupération pondéré & Durée & Temps nécessaire pour que la reprise atteigne 95~\% de son amplitude & Faible \\
SRA, score de robustesse adaptative & Adaptation & Capacité à remonter, pénalisée par un écart durable au niveau initial & Élevé \\
ERR, efficacité de résilience relative & Performance conservée & Part du service nominal maintenue sur l'épisode & Élevé \\
PR, potentiel de récupération & Rétablissement & Niveau final rapporté à la mi-profondeur de la chute & À préciser \\
\bottomrule
\end{tabularx}
\end{table}

Six indicateurs sont calculés analytiquement à partir des paramètres ajustés. Ils héritent ainsi du lissage opéré par l'ajustement et décrivent chacun une facette de la résilience : perte subie, vitesse et durée de la reprise, capacité d'adaptation, performance conservée et potentiel de rétablissement.

\begin{align}
\mathrm{IPC} &= \frac{h}{\alpha}\ln 2 + \frac{h\,d}{2} - \frac{h + q}{\beta}\ln 2 \label{eq:S5} \\
\mathrm{CRD} &= \frac{\beta}{\alpha}\cdot\frac{h + q}{h}\cdot\frac{1}{d}\cdot\sqrt{\alpha\beta} \label{eq:S6} \\
\mathrm{TRP} &= d + \frac{2{,}95}{\beta} \label{eq:S7} \\
\mathrm{SRA} &= \frac{\beta\,(h + q)}{\alpha\,h\,d}\, e^{-|q|/h}\,\big(1 + \tanh(\beta - \alpha)\big) \label{eq:S8} \\
\mathrm{ERR} &= 1 + \frac{q}{2k} - \frac{h\,|\mathrm{IPC}|}{T\,k}, \quad T = g + d + \frac{4}{\beta} \label{eq:S9} \\
\mathrm{PR} &= \frac{k + q}{k - h/2} \label{eq:S10}
\end{align}

L'horizon $T$ de l'équation~\eqref{eq:S9} couvre la chute et une reprise complète à environ 98 \%, conformément à la définition proposée par \citet{rahiel2026in4pl} ; il s'adapte ainsi à la durée propre de chaque crise au lieu d'imposer une fenêtre fixe.

\subsection{Qualité d'ajustement, cas limites et domaine de la résilience}
\label{sec:S3.3}

La qualité de chaque ajustement est appréciée par le coefficient de détermination et l'erreur quadratique moyenne, calculés sur la trajectoire brute, ainsi que par un coefficient de détermination calculé sur la trajectoire lissée sur cinq jours. L'écart entre les deux coefficients distingue un défaut de forme de la courbe d'une simple fluctuation journalière de la demande. Les crises dont l'ajustement est insuffisant, ou dont la chute est trop faible pour être caractérisée, peuvent être écartées de la base.

Les crises absorbées sans dégradation notable, lorsque le réseau dispose d'assez de marge ou, en représentation temporelle, d'assez de stock, posent une première difficulté : la profondeur $h$ tend vers zéro et les indicateurs rapportés à $h$, CRD et SRA, deviennent instables. La couverture de stock permet de déplacer de manière contrôlée le seuil à partir duquel une crise est ainsi absorbée. Les trajectoires à plusieurs creux se rencontrent typiquement lorsqu'un stock mobilisé s'épuise avant la fin de la crise, ou lorsque le retour des capacités provoque un rétablissement brutal : la courbe n'en restitue qu'une enveloppe. Enfin, PR augmente avec la profondeur de la chute à niveau final donné ; son sens d'interprétation doit être arrêté avant de l'utiliser comme objectif de pilotage.

Ces cas limites tiennent à la définition même de la résilience. Trois propriétés voisines se distinguent par la forme de la trajectoire du taux de service. La robustesse est la capacité à absorber un choc sans dégradation notable : le taux de service reste au niveau nominal et il n'y a rien à rétablir. La stabilité est la capacité à rester au voisinage de ce niveau malgré des variations rapides et de faible amplitude, comme les fluctuations journalières de la demande, qui relèvent de la régulation courante. La résilience suppose au contraire une chute effective, suivie d'un rétablissement : elle se mesure à la perte subie et à la manière dont le taux de service revient à un niveau acceptable. Son domaine est donc borné des deux côtés. Sans chute, la question de la résilience ne se pose pas ; lorsque la chute est telle que le réseau cesse durablement de servir et ne se rétablit pas sur la période observée, la crise relève de la défaillance, qui appelle une reconfiguration plutôt qu'un rétablissement. La courbe hyperbolique composée \citep{rahiel2026in4pl} est construite pour ce domaine, qui va d'une chute à une reprise. L'instabilité de CRD et SRA lorsque $h$ tend vers zéro signale ainsi une crise absorbée, c'est-à-dire un cas de robustesse, et non une limite du modèle.

\subsection{Signatures des dysfonctionnements sur la courbe}
\label{sec:S3.4}

Les paramètres de la courbe relient la trajectoire observée aux mécanismes décrits en section 3 du manuscrit. Le tableau~\ref{tab:S3} résume les signatures attendues ; il constitue un ensemble d'hypothèses que l'analyse des résultats et l'apprentissage causal permettent de confirmer ou de nuancer.

\begin{table}[htbp]
\caption{Signatures des dysfonctionnements sur la courbe de résilience.}
\label{tab:S3}
\footnotesize
\begin{tabularx}{\textwidth}{@{}LLL@{}}
\toprule
\textbf{Phénomène} & \textbf{Signature attendue sur la courbe} & \textbf{Origine probable} \\
\midrule
Absorption du choc & $h$ faible, $g$ élevé & Marges de capacité, voies de secours suffisantes \\
Absorption par les stocks & $h$ faible ou nul, $g$ élevé & Couverture de stock élevée au regard de la durée du choc \\
Rupture brutale & $h$ élevé, $\alpha$ élevé & Coupure ou fermeture sévère, réseau sous tension \\
Rupture différée & $g$ élevé, puis $h$ élevé & Stocks épuisés en cours de crise, matière en route insuffisante \\
Dégradation progressive & $\alpha$ faible & Congestion diffuse, pic de demande modéré \\
Crise prolongée & $d$ élevé & Perturbation longue, détection tardive, inertie décisionnelle \\
Reprise retardée & $d$ supérieur à la durée du choc & Reconstitution des stocks et de la matière en route après le retour des capacités \\
Reprise rapide & $\beta$ élevé, TRP faible & Levier efficace engagé tôt \\
Perte durable & $q$ $<$ 0 & Capacité non restaurée à la fin de la période, stock épuisé \\
Sur-récupération & $q$ $>$ 0 & Leviers maintenus après la fin de la perturbation \\
\bottomrule
\end{tabularx}
\end{table}
